\section{Continuous source resolution and abelian reconstruction}\label{sec:profinite-reconstruction}

This section keeps the profinite source fixed while allowing its finite images to grow. It separates two independent resolution coordinates: coefficient precision and the finite quotient through which a particular observation factors. Its strongest constructive conclusion is a complete finite-matrix lifting criterion for continuous rank-two representations of $\mathbb Z_p^r$. In arbitrary rank the same criterion holds when a specified normalized generator is cyclic modulo $p$. No prescribed determinant, ramification type, Hodge condition, or automorphic realization is included in these statements.

The continuous cohomology used below is the cohomology of continuous inhomogeneous cochains with a \emph{finite discrete} coefficient module. The finite-quotient description is classical \cite[Chapter II, Section 4, Proposition 4.2]{MilneCFT2020}; see also \cite{StacksContinuous}. The free pro-$p$ construction and cyclic-centralizer argument are given directly. The claims are written deductions under explicit hypotheses, not assertions of historical novelty.

\subsection{The actual continuous obstruction}

Fix a profinite group $\Gamma$, a prime $p$, $n,d\geq1$, and a continuous representation
\[
 \rho_n:\Gamma\longrightarrow\operatorname{GL}_d(R_n),
 \qquad R_n=\mathbb Z/p^n.
\]
The target is finite and discrete. Give $V=M_d(\mathbb F_p)$ the adjoint action of $\bar\rho$. Choose a section of the reduction map on the finite set $\rho_n(\Gamma)$, taking the identity to the identity. Composing this section with $\rho_n$ gives a continuous matrix-valued lift $A_g$ to $R_{n+1}$, not necessarily multiplicative. Define
\[
 c_A(g,h)=\frac{A_gA_hA_{gh}^{-1}-I}{p^n}\pmod p.
\]
Here division is the identification of $p^nR_{n+1}$ with $\mathbb F_p$, as in \cref{thm:relation-H2}.

\begin{theorem}[Complete continuous one-step obstruction]\label{thm:continuous-H2}
The class
\[
 o_n^{\mathrm{cts}}(\rho_n)=[c_A]\in H^2_{\mathrm{cts}}(\Gamma,V)
\]
is independent of the continuous choices. It vanishes if and only if $\rho_n$ has a \emph{continuous} lift to $\operatorname{GL}_d(R_{n+1})$ with its specified reduction. More precisely, the repair equation is
\[
 \delta b=-c_A,\qquad b:\Gamma\to V\text{ continuous and normalized},
\]
and the repaired matrices are $(I+p^nb(g))A_g$. When lifts exist, their framed set is a torsor for $Z^1_{\mathrm{cts}}(\Gamma,V)$. Strict-conjugacy classes and stabilizers are governed by $H^1_{\mathrm{cts}}$ and $H^0_{\mathrm{cts}}$, respectively.
\end{theorem}
\begin{proof}
The associativity and square-zero computations in \cref{thm:relation-H2} hold pointwise without a finiteness assumption on the source. All the resulting cochains are continuous because the targets are finite and the chosen lifts are continuous. Any other continuous lift differs uniquely by $(I+p^nb(g))$ with continuous $b$, and its residual is $c_A+\delta b$. Hence a zero class is equivalent to an actual continuous repair. Conversely an actual continuous homomorphism is such a repair. The difference of two repairs is a continuous $1$-cocycle. Conjugation by $I+p^nX$ changes it by a coboundary, with kernel the invariant elements of $V$, exactly as in the finite case.
\end{proof}

\begin{theorem}[Source refinement and inflation]\label{thm:source-inflation}
Choose an open normal $U\subseteq\ker\rho_n$ and compute $c_A$ on $Q=\Gamma/U$. Then
\[
 o_n^{\mathrm{cts}}(\rho_n)=0
 \quad\Longleftrightarrow\quad
 \exists\, U'\triangleleft_{\mathrm{open}}\Gamma,\ U'\subseteq U:
 \operatorname{Inf}_{\Gamma/U}^{\Gamma/U'}[c_A]=0.
\]
Thus a continuous next-level lift exists exactly when a repair exists at \emph{some finer source quotient}, not necessarily the quotient used initially. For every $j\geq0$,
\[
 H^j_{\mathrm{cts}}(\Gamma,V)
 \cong\underset{U'\subseteq U}{\operatorname{colim}}
 H^j(\Gamma/U',V).
\]
\end{theorem}
\begin{proof}
Every continuous map $\Gamma^j\to V$ is locally constant. Cover its compact domain by finitely many products of cosets on which the map is constant. Intersect the finitely many open normal subgroups defining these cosets and further intersect with $U$. The map then factors through $(\Gamma/U')^j$. This applies to cocycles, coboundaries, and cochains witnessing an equality. The finite-quotient cochain complexes therefore have the continuous complex as their filtered union, and the same observation about witnesses proves the asserted cohomology colimit. In particular a class vanishes in that colimit exactly when it vanishes at some subsequent stage. Apply \cref{thm:continuous-H2}. This is the finite-module specialization of the cited classical inflation theorem.
\end{proof}

\begin{example}[An explicit obstruction killed by refining the source]\label{ex:inflation-carry}
Let $p$ be odd, $J\in M_d(\mathbb Z)$ nonzero modulo $p$, and $A=I+pJ$ modulo $p^3$. Then $A^p=I+p^2J$ and $A^{p^2}=I$ modulo $p^3$. The representation of $C_p$ modulo $p^2$ uses powers of $A$. For representatives $0\leq a,b<p$, choose $A_a=A^a$ at the next level. Its cocycle is
\[
 c(a,b)=\left\lfloor\frac{a+b}{p}\right\rfloor J\pmod p.
\]
The adjoint residual action is trivial. For any normalized cochain $u$ on $C_p$,
\[
 \sum_{t=0}^{p-1}(\delta u)(1,t)=0,
 \qquad\sum_{t=0}^{p-1}c(1,t)=J\ne0.
\]
So the finite-quotient class is nonzero. On $C_{p^2}$ put
\[
 b(x)=\left\lfloor\frac{x}{p}\right\rfloor J\pmod p
 \quad(0\leq x<p^2).
\]
Writing each $x$ as its low digit plus $p$ times its high digit shows
$\delta b=-\operatorname{Inf}(c)$: the high-digit carry is the negative of the low-digit carry modulo $p$. The repaired map is exactly $x\mapsto A^x$. Thus inflation kills a genuine finite-quotient obstruction. This calculation does not declare every continuous obstruction removable.
\end{example}

\subsection{Reconstructing a continuous tower}

\begin{theorem}[Continuous reconstruction from unobstructed stages]\label{thm:continuous-tower}
If $H^2_{\mathrm{cts}}(\Gamma,\operatorname{Ad}\bar\rho)=0$, every given $\rho_n$ with that residual action extends to a continuous representation
\[
 \rho:\Gamma\longrightarrow\operatorname{GL}_d(\mathbb Z_p).
\]
Without the cohomology-vanishing hypothesis, a compatible sequence of continuous lifts still reconstructs uniquely as a framed representation. If $\Gamma$ is topologically finitely generated, existence of \emph{some} extension of the fixed $\rho_n$ to every coefficient level is enough to obtain a compatible sequence and hence an integral lift.
\end{theorem}
\begin{proof}
The residual adjoint module is unchanged by each coefficient lift. Under the vanishing hypothesis, \cref{thm:continuous-H2} permits recursive choice of compatible lifts $\rho_m$ for all $m\geq n$. Take the inverse limit of their matrix entries. Multiplication is compatible with this limit, and determinants remain units because their reductions modulo $p$ are units. Each finite projection is continuous, so the map into the inverse-limit topology is continuous. The projections determine it uniquely.

For the last assertion, choose $r$ topological generators. A continuous map into a finite group is determined by their images, since the equalizer of two such maps is closed and contains a dense subgroup when the images agree. Thus there are finitely many continuous representations at each level. Form the rooted tree of extensions of $\rho_n$, with edges given by one-step reduction. It is finitely branching. If every level is nonempty, the root has descendants at arbitrarily large depth. At least one child has the same property, since there are only finitely many children. Repeating gives an infinite compatible path. Apply the first limit argument. This compactness argument supplies no numerical stopping bound and no algorithm for computing $H^2_{\mathrm{cts}}$ of an arbitrary source.
\end{proof}

\begin{theorem}[Unrestricted coefficient lifting for a free pro-$p$ source]\label{thm:free-pro-p-lift}
Let $F_r^{(p)}$ be the free pro-$p$ group on finitely many generators. Every continuous
\[
 F_r^{(p)}\longrightarrow\operatorname{GL}_d(R_n)
\]
has a continuous lift to $\operatorname{GL}_d(\mathbb Z_p)$. One may choose arbitrary integral lifts of its generator matrices.
\end{theorem}
\begin{proof}
Construct $F_r^{(p)}$ as the inverse limit of the finite $p$-group quotients of the abstract free group on these generators. A map from the generators to a pro-$p$ group induces compatible maps to each of its finite $p$-group quotients, and hence a unique continuous map from this inverse limit. This proves the universal property needed here.

The residual image $P\subseteq\operatorname{GL}_d(\mathbb F_p)$ is a $p$-group, being a finite quotient of a pro-$p$ group. Let $H$ be its inverse image in $\operatorname{GL}_d(\mathbb Z_p)$. Its quotient modulo $p^m$ is a finite $p$-group: its kernel over $P$ has the filtration by congruence subgroups with additive $M_d(\mathbb F_p)$ subquotients. Therefore $H$ is pro-$p$, including at $p=2$. Arbitrary integral lifts of the generator matrices lie in $H$. The just-proved universal property extends them to a continuous map. Reduction agrees with the input on the generators and therefore everywhere. This does not assert that the lift is faithful or has finite image.
\end{proof}

\subsection{Continuous powers and the cyclic-centralizer mechanism}

\begin{lemma}[When an integral matrix has continuous procyclic powers]\label{lem:procyclic-powers}
Let $A\in\operatorname{GL}_d(\mathbb Z_p)$. The map $m\mapsto A^m$ on $\mathbb Z$ extends uniquely to a continuous homomorphism $\mathbb Z_p\to\operatorname{GL}_d(\mathbb Z_p)$ if and only if the reduction of $A$ modulo $p$ has $p$-power order. In particular a commuting tuple of such matrices determines a continuous representation of $\mathbb Z_p^r$.
\end{lemma}
\begin{proof}
Necessity follows because a continuous image of $\mathbb Z_p$ in a finite group is a cyclic $p$-group. Conversely suppose $A^{p^t}\in I+pM_d(\mathbb Z_p)$. For $B\in I+p^sM_d(\mathbb Z_p)$ with $s\geq1$, binomial expansion shows $B^p\in I+p^{s+1}M_d(\mathbb Z_p)$. For $p=2$ the quadratic term has depth $2s\geq s+1$, so this estimate remains valid. It follows that
\[
 A^{p^{t+j}}\in I+p^{j+1}M_d(\mathbb Z_p).
\]
Hence exponents congruent modulo $p^{t+j}$ give powers congruent modulo $p^{j+1}$, also for negative integer exponents by invertibility. Completeness extends the map continuously from the dense subgroup $\mathbb Z$ to $\mathbb Z_p$. The group law and uniqueness follow by continuity. For commuting matrices the product of these powers defines the final assertion.
\end{proof}

\begin{lemma}[Cyclic centralizers over a coefficient ring]\label{lem:cyclic-centralizer}
Let $R$ be a commutative ring, $B\in M_d(R)$, and suppose a vector $v$ makes
\[
 Q=(v,Bv,\ldots,B^{d-1}v)\in\operatorname{GL}_d(R).
\]
Then every matrix commuting with $B$ is uniquely $f(B)$ for a polynomial $f\in R[t]$ of degree less than $d$. The coefficients of $f$ are the entries of $Q^{-1}Mv$. In particular, the centralizer is commutative.

Over $R=\mathbb Z/p^s$, it is enough that $\bar B$ have a cyclic vector over $\mathbb F_p$. Every nonscalar $2\times2$ matrix over a field has such a vector.
\end{lemma}
\begin{proof}
There is a unique polynomial of degree less than $d$ satisfying $f(B)v=Mv$, since the columns of $Q$ are a basis. If $MB=BM$, then $M-f(B)$ annihilates every $B^jv$, so it is zero. Conversely every polynomial in $B$ commutes with $B$. This also proves uniqueness and the coefficient formula. A lifted cyclic vector modulo $p$ makes $\det Q$ a unit modulo $p^s$.

Finally, if a $2\times2$ matrix has no cyclic vector, it takes each vector into its own span. On a basis $e_1,e_2$ it is diagonal, and applying this condition to $e_1+e_2$ makes the two diagonal entries equal. Thus it is scalar. This argument includes the field $\mathbb F_2$.
\end{proof}

\subsection{A complete normalized-commutator obstruction}

Here the source is $\Gamma=\mathbb Z_p^r$, where $r\geq1$ is finite. Write the generator matrices of the given continuous representation $\rho_n$ as $A_1,\ldots,A_r$. They commute in $M_d(R_n)$ and each has $p$-power order modulo $p$.

Choose $0\leq a<n$ and integer scalars $\lambda_i$ such that
\[
 A_i\equiv\lambda_i I\pmod{p^a}\quad\text{for every }i.
\]
There is no restriction in this display when $a=0$. Define the \emph{normalized matrices}
\[
 B_i=\frac{A_i-\lambda_i I}{p^a}
 \quad\in M_d(R_{n-a}).
\]
This is well-defined from the given finite data. It does not invert $p$ inside $R_n$. Assume that, for some distinguished index $*$, the matrix $B_*$ is cyclic modulo $p$. Define the finite obstruction record
\[
 \boxed{\Omega_{a,*}(\rho_n)=
 \bigl([B_*,B_i]\bigr)_{i=1}^r\in M_d(R_{n-a})^r,}
 \qquad [B,C]=BC-CB.
\]

\begin{theorem}[Normalized cyclic lifting criterion]\label{thm:normalized-cyclic-lift}
Under the preceding hypotheses,
\[
 \boxed{\rho_n\text{ has a continuous integral lift}
 \quad\Longleftrightarrow\quad\Omega_{a,*}(\rho_n)=0.}
\]
If the obstruction vanishes, an integral lift is constructed from finite polynomial certificates. Express $B_i=f_i(B_*)$ in $R_{n-a}$, with $\deg f_i<d$ and $f_*(t)=t$ when $d>1$. Choose an integral lift $\widehat B$ of $B_*$ and integer lifts $\widehat f_i$ of these polynomials, taking the distinguished generator itself to $\widehat B$. Then put
\[
 \widehat A_i=\lambda_i I+p^a\widehat f_i(\widehat B),
 \qquad
 \rho(x_1,\ldots,x_r)=\prod_i\widehat A_i^{x_i}.
\]
For $d=1$, use the scalar matrices directly. The construction preserves the prescribed reduction modulo $p^n$ and does not pass to an obstruction quotient.
\end{theorem}
\begin{proof}
Suppose first that $\widehat A_i$ are any commuting integral lifts. Then
$\widehat B_i=(\widehat A_i-\lambda_iI)/p^a$ are integral. The identities
\[
 0=[\widehat A_*,\widehat A_i]
   =p^{2a}[\widehat B_*,\widehat B_i]
\]
imply $[\widehat B_*,\widehat B_i]=0$, since matrices over $\mathbb Z_p$ have no $p$-torsion. Reducing modulo $p^{n-a}$ proves necessity. This argument permits arbitrary higher digits in all the integral lifts.

Conversely, the cyclic-centralizer lemma supplies the unique polynomials $f_i$ from the vanishing commutators. Their integral polynomial evaluations commute exactly. Multiplication by $p^a$ shows that $\widehat A_i$ reduces to $A_i$ modulo $p^n$. Each is invertible because its reduction modulo $p$ is invertible. That reduction has $p$-power order by the input's continuity. The procyclic-powers lemma now supplies the continuous representation, whose reduction agrees on the topological generators and hence on all of $\mathbb Z_p^r$. In rank one the same argument is simply scalar lifting. Polynomial coefficients and a cyclic vector give a finite positive certificate; a nonzero entry of a displayed commutator gives a finite negative certificate under the verified cyclic hypothesis.
\end{proof}

\begin{remark}[What normalization reveals]
The input relations only force
\[
 p^{2a}[B_i,B_j]\equiv0\pmod{p^n}.
\]
When $a>0$, this can conceal a nonzero commutator in $R_{n-a}$. The obstruction is therefore not just the original zero commutator re-labelled. It measures the additional commutativity required after removing a common scalar precision. Different permitted scalar choices change $B_i$ by scalar matrices and leave its commutators unchanged. Changing representatives changes them only by zero in $R_{n-a}$. No preferred normalized matrix is claimed independently of these data.
\end{remark}

\begin{corollary}[A source-fixed finite horizon in the cyclic stratum]\label{cor:abelian-cyclic-horizon}
Under the hypotheses of \cref{thm:normalized-cyclic-lift}, the following are equivalent:
\[
 \rho_n\text{ lifts over }\mathbb Z_p,
 \qquad \rho_n\text{ extends over }R_{n+a},
 \qquad \Omega_{a,*}(\rho_n)=0.
\]
The extension remains a representation of the same $\mathbb Z_p^r$. When $a=0$, the current representation is already a sufficient certificate.
\end{corollary}
\begin{proof}
An integral lift supplies the finite extension. Given a finite extension with commuting matrices $A_i'$ modulo $p^{n+a}$, subtract the same $\lambda_iI$ and divide by $p^a$. The resulting matrices $B_i'$ are defined modulo $p^n$ and reduce to $B_i$ modulo $p^{n-a}$. The equality
$p^{2a}[B_*',B_i']=0$ modulo $p^{n+a}$ implies
$[B_*',B_i']=0$ modulo $p^{n-a}$, so $\Omega_{a,*}=0$. The normalized criterion supplies an integral lift. This is a bound obtained from the source's actual commutativity relations, not from the orders of its growing finite images.
\end{proof}

\subsection{Rank two: a complete decision from the given precision}

For $A\in M_2(R_n)$ define its scalar depth by
\[
 a(A)=\max\{0\leq t\leq n:A\text{ is scalar modulo }p^t\}.
\]
Writing entries in $R_n$, this is the minimum of the capped valuations of
$A_{12},A_{21},A_{11}-A_{22}$. It is unchanged by invertible conjugation. For a tuple put $a=\min_i a(A_i)$.

\begin{theorem}[Complete rank-two profinite lifting test]\label{thm:rank-two-profinite}
Let $p$ be any prime and let $\rho_n:\mathbb Z_p^r\to\operatorname{GL}_2(R_n)$ be continuous, with finitely many topological generators as above.
\begin{enumerate}
\item If $a=n$, all the $A_i$ are scalar and $\rho_n$ has a continuous integral lift.
\item If $a<n$, choose any generator $A_*$ with scalar depth $a$. Take $\lambda_i$ to be an integer representative of $(A_i)_{11}$ and form $B_i$ modulo $p^{n-a}$. Then
\[
 \boxed{\rho_n\text{ lifts over }\mathbb Z_p
 \quad\Longleftrightarrow\quad
 [B_*,B_i]=0\text{ in }M_2(R_{n-a})\ (1\leq i\leq r).}
\]
The test is finite matrix arithmetic on the \emph{existing input}. It does not require enumeration of deeper coefficient lifts. When it passes, the polynomial construction in \cref{thm:normalized-cyclic-lift} produces exact commuting integral generator matrices.
\end{enumerate}
\end{theorem}
\begin{proof}
If all generators are scalar, lift their scalar entries integrally. Their residues modulo $p$ have $p$-power order and therefore equal $1$, since $\mathbb F_p^\times$ has order prime to $p$. The scalar lifts are units and define continuous powers.

Otherwise every $A_i$ is scalar modulo $p^a$, while $B_*$ is nonscalar modulo $p$ by the definition of $a$. It is cyclic by \cref{lem:cyclic-centralizer}. Thus every hypothesis of \cref{thm:normalized-cyclic-lift} holds. In rank two the polynomials have the form $u_i+v_it$ and can be recovered from a two-column cyclic basis by solving one invertible $2\times2$ linear system modulo $p^{n-a}$. The equivalence with liftability also shows independence of the distinguished generator and of basis choices, although the selected lift itself need not be canonical.
\end{proof}

\begin{corollary}[A sharp uniform rank-two horizon]\label{cor:rank-two-sharp-horizon}
For every $p,n,r$ as above,
\[
 \boxed{\rho_n\text{ lifts continuously over }\mathbb Z_p
 \quad\Longleftrightarrow\quad
 \rho_n\text{ extends continuously to }R_{2n-1}.}
\]
For $n\geq2$ and $r\geq2$, $2n-1$ is optimal as a uniform coefficient level. For $n=1$, all continuous rank-two representations of $\mathbb Z_p^r$ over $\mathbb F_p$ lift.
\end{corollary}
\begin{proof}
For a nonscalar tuple, $a\leq n-1$, so an extension to $2n-1$ restricts to one at $n+a$. Apply \cref{cor:abelian-cyclic-horizon}. Scalar tuples lift directly. This also proves the assertion at $n=1$.

For sharpness fix $n\geq2$ and put
\[
 A=I+p^{n-1}E_{12},\qquad B=I+p^{n-1}E_{21}.
\]
These matrices commute modulo $p^{2n-2}$ and generate a finite abelian $p$-group there, so they give a continuous extension of their level-$n$ representation to level $2n-2$. However every pair of lifts preserving level $n$ has the form
$\widetilde A=A+p^nP$, $\widetilde B=B+p^nQ$, and
\[
 \widetilde A\widetilde B-\widetilde B\widetilde A
 \equiv p^{2n-2}
 \begin{pmatrix}1&0\\0&-1\end{pmatrix}
 \not\equiv0\pmod{p^{2n-1}}.
\]
Indeed every correction term has valuation at least $2n-1$. This proves nonextension at $2n-1$ for every prime, including $2$, and hence nonexistence of an integral lift. Add identity generators for $r>2$. No smaller uniform coefficient level can distinguish these inputs from extendable ones by mere existence of an extension.
\end{proof}

\begin{example}[The first hidden commutator]\label{ex:abelian-tangent}
Suppose $n=2$ and all generator matrices are $A_i=I+pX_i$ modulo $p^2$. They always commute at that coefficient level. For rank two, \cref{thm:rank-two-profinite} gives
\[
 \rho_2\text{ lifts over }\mathbb Z_p
 \quad\Longleftrightarrow\quad
 [X_i,X_j]=0\text{ in }M_2(\mathbb F_p)\text{ for all }i,j.
\]
If some $X_i$ is nonscalar, its centralizer is commutative and the displayed test agrees with the distinguished-generator criterion; if all are scalar the lift is immediate. For $X_1=E_{12}$ and $X_2=E_{21}$ it fails. Each generator separately has a continuous integral procyclic lift, but the pair does not define an integral action of $\mathbb Z_p^2$. On the free pro-$p$ source the same pair \emph{does} lift, by \cref{thm:free-pro-p-lift}. The surviving obstruction is therefore an actual source relation, not finite-order torsion imposed by an incidental quotient.
\end{example}

\begin{example}[A positive repair that changes only unresolved digits]\label{ex:abelian-positive-repair}
Work at $p=3$, $n=3$. Take
\[
 A_1=\begin{pmatrix}1&3\\9&1\end{pmatrix},\qquad
 A_2=\begin{pmatrix}7&12\\9&7\end{pmatrix}\pmod{27}.
\]
They commute modulo $27$ and both are the identity modulo $3$. Choose $a=1$, $\lambda_1=\lambda_2=1$. Then
\[
 B_1=\begin{pmatrix}0&1\\3&0\end{pmatrix},\qquad
 B_2=\begin{pmatrix}2&4\\3&2\end{pmatrix}\pmod9,
 \qquad B_2=2I+4B_1\pmod9.
\]
The vector $e_2$ is cyclic for $B_1$ modulo $3$. The displayed integer representatives of $A_1,A_2$ do not commute exactly: their additive commutator is $\operatorname{diag}(-81,81)$. The criterion nevertheless passes. Retain $B_1$ integrally and replace $B_2$ by the exact polynomial $2I+4B_1$. This constructs
\[
 \widehat A_1=\begin{pmatrix}1&3\\9&1\end{pmatrix},\qquad
 \widehat A_2=\begin{pmatrix}7&12\\36&7\end{pmatrix}.
\]
They commute exactly, remain units over $\mathbb Z_3$, and have the original reduction modulo $27$. Their continuous powers give the integral representation. This is an explicit non-destructive repair of a tuple, not a claim that its initial arbitrary integer representatives were already an exact action.
\end{example}

\subsection{The repaired scope of the two-axis calculus}

There is a necessary variance distinction. Let $\mathcal Q_\Gamma$ have open normal subgroups as objects and an arrow $U\to U'$ when $U'\subseteq U$. Write $\mathscr P_d(n,U)$ for the groupoid of rank-$d$ representations of $\Gamma/U$ on $R_n^d$, with invertible intertwining matrices as arrows.

\begin{proposition}[The source--coefficient diagram has mixed variance]\label{prop:mixed-source-variance}
Reduction of coefficients and inflation of source define a functor
\[
 \mathscr P_d:\mathbb N_{\geq1}^{op}\times\mathcal Q_\Gamma
 \longrightarrow\mathbf{Gpd}.
\]
Each source-inflation functor is fully faithful, and it commutes with coefficient reduction. An object at $(m,U')$ is a compatible refinement of $\rho_n$ at $(n,U)$ precisely when its reduction to level $n$ equals, or is supplied with an isomorphism to, the inflation of $\rho_n$. There is no total source-coarsening functor on arbitrary representations obtained simply by reversing inflation.
\end{proposition}
\begin{proof}
Coefficient reduction acts on the matrix entries of representations and intertwiners. Source inflation precomposes them with the surjection $\Gamma/U'\to\Gamma/U$. The two operations commute directly on both objects and arrows. An intertwiner of inflated representations commutes with the same matrix image as before inflation, because the source map is onto. Thus inflation is fully faithful. Finally a representation of $\Gamma/U'$ descends through $\Gamma/U$ only if the subgroup $U/U'$ acts trivially. This fails for general representations, as the explicit repaired action in \cref{ex:inflation-carry} shows at its higher coefficient level.
\end{proof}

Coefficient refinement and source refinement do different jobs. The former improves arithmetic accuracy. The latter permits more of the source to act, and can kill an inflated finite-quotient obstruction without deleting coefficient information. Neither permits violating a relation of the fixed source itself. In particular these two resolution coordinates cannot both be inserted into the earlier coarsening-only convention without recording their opposite variance or restricting to descending objects.

The constructive abelian theorems avoid an unbounded quotient search altogether on their stated strata. Scalar-depth normalization reveals the missing relation, the cyclic centralizer supplies a finite polynomial certificate, and continuous matrix powers reconstruct the action. In rank two this yields a complete test at every coefficient precision and every finite source rank. In larger rank it requires the stated cyclic normalized generator. Free pro-$p$ sources furnish a contrasting unrestricted coefficient-lifting case because they impose no additional relations on the generators.

These results do not imply a universal scalar-depth test in higher rank, an unconstrained theorem for all profinite groups, or a lift preserving extra arithmetic deformation conditions. The constructions make choices of higher matrix coefficients and may change determinants or other invariants not included among the equations. Continuous Galois representations with local geometric conditions and automorphic existence remain separate targets. The next substantive extension is a specified noncommutative profinite relation system with controlled centralizers or a verified deformation complex, rather than an assumption that every vanishing selected diagnostic gives a global object.
